\chapter{Visitors}
\label{ch:tourism}

\section{The industry Cuba built, and what it cost}

Cuba rebuilt tourism from almost nothing after 1990 --- 340,000 visitors in that
year, 4.7 million in 2018 --- and by any conventional measure that is a success.
The trouble is what happened next, and it is the clearest case study in this book
of investment allocated by power rather than by return.

Through the decade to 2024 the Cuban state, principally through the
military-linked conglomerate GAESA and its hotel arm, built hotel rooms at a
sustained and increasing rate while occupancy collapsed.

The 2024 figures are worth setting out in full, and in shares rather than in
pesos, because shares are immune to the exchange-rate problem of the note on
sources. Of total national investment that year, business services, real estate
and rentals took 24.9 billion pesos and hotels and restaurants a further 11.9
billion --- together 37.4 per cent of everything the country invested.
Agriculture, livestock and forestry took 2.6 billion, or 2.7 per cent. Public
health and social assistance took 2.0 billion, less than a fifth of what went
into hotels and restaurants alone. As the economist Pedro Monreal put it,
agricultural investment was fourteen times smaller than tourism's
\citep{oneiinv}.\unv{}

That is the allocation of a country importing most of its food, unable to light
its cities, appealing to the World Food Programme for children's milk, and
running its existing hotels at occupancy rates reported around a quarter. New
hotels opened into a market with no guests while the pumping stations that
supply those same hotels with water failed for want of parts.

No market signal produced that, and no market signal could have. It was
produced by an institution that could
allocate capital to itself, without publishing a return, without competing for
the funds, and without anyone in a position to say no. Chapter~\ref{ch:integrity}
uses this as its worked example, because it demonstrates that Cuba's investment
problem is not primarily a shortage of capital. It is that the capital which
exists is allocated by whoever controls it, to whatever they control.

Two other features of the industry as built matter for what follows.

It was an \emph{enclave}. The dominant form is the all-inclusive beach resort at
Varadero and the northern cays, physically separated from the country around it,
supplied largely with imported food, staffed by workers paid in pesos through a
state agency that keeps the currency spread, and --- until 2008 --- legally
closed to Cubans. The Caribbean-wide evidence on this model is that the share of
the visitor's spending which stays in the host country is low, commonly cited in
the range of a fifth to a third for all-inclusives, with the remainder accruing
to the foreign tour operator, the foreign airline and the imported supply
chain.\unv{} Cuba adopted the model that leaks most.

And it carries a memory. Chapter~\ref{ch:inheritance} described the Havana of
the 1950s: casinos, organised crime, a pleasure district for foreigners in a
country whose own people were excluded. The revolution's founding claim included
abolishing that. Any proposal to make tourism an anchor industry has to answer
the objection that it recreates it, and the answer cannot be that times have
changed. It has to be a design in which Cubans own the businesses, are not
excluded from the spaces, and capture the income.

\section{The strategy: value per visitor, not visitors}

The book proposes an explicit reversal of the metric the industry has been run
on.

Cuba should not pursue arrival numbers. It should pursue \emph{net foreign
exchange retained per visitor}, and should be willing to accept fewer visitors
to raise it.

The reasoning is partly environmental and partly arithmetic. Chapter~\ref{ch:nature}
sets out what mass tourism does to reefs, wetlands and colonial fabric, and Cuba's
comparative advantage is precisely the thing that volume destroys. But the
arithmetic is sufficient on its own: a visitor in a private room eating in
private restaurants and buying from private guides leaves several times as much
in Cuban hands as a visitor in an all-inclusive, even when the second one paid
more for the holiday. On a labour-constrained island with a failing grid, adding
visitors is expensive --- each one needs water, power, food and staff, all of
which are scarce --- while adding value per visitor is nearly free.

This also resolves the tension with Chapter~\ref{ch:premises}'s labour
constraint. Cuba does not have the people for a Cancún. It does have the people
for something better paid.

\begin{design}{The visitor economy}
\label{d:visitor}
\begin{rules}
\rl{1}{The metric} The published objective of tourism policy is net foreign
exchange retained per visitor-night, reported annually with its components. No
arrival target is set or published. What a state measures is what it builds, and
Cuba built rooms because it counted arrivals.

\rl{2}{Private by default} Accommodation, food service, guiding, transport and
retail to visitors are open to private and cooperative providers by right, under
Design~\ref{d:commit}.2, with the state's role confined to standards, safety and
inspection. The \emph{casa particular} and the \emph{paladar} are the default
form of the Cuban tourism industry, not its tolerated margin.

\rl{3}{Direct employment and direct payment} The state employment agency
arrangement is abolished for this sector as for all others
(Chapter~\ref{ch:capital}): hotels hire their own staff and pay them directly.
The currency spread taken from tourism workers is the industry's founding
injustice and it is also why the service is indifferent.

\rl{4}{Local procurement} Accommodation businesses above a size threshold must
publish the share of food purchased domestically, and the accommodation levy of
Design~\ref{d:tax}.7 is set at two rates, the lower applying above a published
domestic-procurement threshold. This is the instrument that connects
Chapter~\ref{ch:land} to this chapter: Cuban agriculture's most reliable
potential customer is standing in a hotel two hours away, buying imported
tomatoes.

\rl{5}{No new state hotels} A moratorium on state hotel construction until
system-wide occupancy exceeds a published threshold for two consecutive years.
The capital released is redirected under Chapter~\ref{ch:sequence}'s financing
table. This is the single largest reallocation of investment available to Cuba
and it requires no foreign money at all.

\rl{6}{Open to Cubans} No facility, beach, cay or district is closed to Cuban
residents, in law or in practice, and the prohibition is justiciable under
Design~\ref{d:courts}.4.
\end{rules}
\end{design}

\section{What Cuba actually has}

The asset is not beach. Every country in the Caribbean has beach, several have
better, and competing on it means competing on price with places that have
better airports and working electricity.

What Cuba has that the competition does not is a combination that cannot be
built: nine UNESCO World Heritage sites; four intact colonial cities; a musical
culture that is a destination in itself; some of the best-preserved coral in the
hemisphere at the Jardines de la Reina; the largest wetland in the insular
Caribbean at Zapata, with the endemic bird list that goes with it; five thousand
seven hundred kilometres of coast, much of it undeveloped; and --- unusually and
importantly --- a very low violent crime rate, which is a decisive advantage in a
region where several competitors have the opposite reputation.

The products that follow from that list are diving, birding and nature tourism,
cultural and musical tourism, architectural and heritage tourism, medical
tourism under Chapter~\ref{ch:life}, and long-stay visitors under
Chapter~\ref{ch:talent} --- all of which are higher-value, lower-volume, less
seasonal and less environmentally destructive than beach resorts, and all of
which are currently underdeveloped or absent.

\subsection{The Old Havana precedent}

Cuba has already built the best heritage institution in the Caribbean and then
dismantled it, and the episode is instructive in both directions.

The Office of the City Historian of Havana, under Eusebio Leal, was given in 1993
something no other Cuban body had: the right to run commercial enterprises in
Old Havana --- hotels, restaurants, shops --- and to \emph{retain and reinvest its
own revenue} in restoring the district. It worked. Old Havana's restoration is
the most successful physical project of the Cuban period, it was self-financing,
it employed and trained local residents, it restored housing alongside monuments
rather than displacing people, and it produced a district that is now the single
strongest tourism asset in the country.

In 2016 its commercial arm was transferred to the military conglomerate, and the
Office's independent revenue base with it.

Both halves belong in the design. The mechanism --- a heritage authority that
keeps its own earnings and is accountable for the fabric --- is correct and
should be extended to Trinidad, Cienfuegos, Camag\"uey and Santiago. And the
reversal is the reason it has to be entrenched under Design~\ref{d:commit}
rather than granted, because it was granted, it worked for twenty-three years,
and it was taken away.

\begin{design}{Heritage authorities}
\label{d:heritage}
\begin{rules}
\rl{1}{Form} A statutory heritage authority for each historic city, with
jurisdiction over planning consent, building standards and public space in its
district.

\rl{2}{Revenue} It retains a defined share of the accommodation levy and of
concession income raised in its district, by statute and not by grant, and this
share may be altered only by legislation.

\rl{3}{Obligation} A published, audited restoration programme, and a binding
requirement that a fixed share of restored floorspace remains residential and
that existing residents have a right to rehousing within the district. Heritage
tourism that empties a district of its inhabitants produces a stage set and
loses the thing visitors came for.

\rl{4}{Independence} Its commercial arm may not be transferred to any other
entity without legislation. Written because this is exactly what happened.
\end{rules}
\end{design}

\section{Carrying capacity}

Chapter~\ref{ch:nature} makes the environmental argument at length. What belongs
here is that it is an economic argument, not a sentimental one.

The Cuban reefs, the Zapata wetland and the cays are the physical stock that the
high-value products above are sold from. A reef that is bleached, dynamited by
anchors or silted by coastal construction stops producing revenue permanently,
and the revenue it was producing was the highest per visitor in the industry.
Diving tourism at the Jardines de la Reina works precisely because the fishing
was excluded and the sharks are still there; it is worth more per visitor-day
than any beach resort in the country and it exists only because someone drew a
line on a map and enforced it.

\begin{design}{Limits}
\label{d:limits}
\begin{rules}
\rl{1}{Site limits} Published daily and annual visitor limits for reef, cay and
protected-area sites, set by the environmental regulator of
Chapter~\ref{ch:nature} and not by the tourism ministry, with access allocated by
published auction or ballot rather than by administrative permission.

\rl{2}{Coastal setback} A statutory construction setback from the shoreline,
applied without exception, including to existing state and military-linked
operators. Chapter~\ref{ch:nature} sets the distance; this clause is about the
absence of exceptions, which is where every Caribbean setback rule has failed.

\rl{3}{No new cay causeways} The \emph{pedraplenes} that connect the northern
cays to the mainland have measurably damaged water circulation and the
ecosystems behind them. No further ones are built, and remediation of the
existing ones is assessed.

\rl{4}{Price the scarcity} Where a site is at its limit, the access price rises.
This is the mechanism that converts an environmental constraint into revenue
under Design~\ref{d:visitor}.1 rather than into a queue.
\end{rules}
\end{design}

\section{Dutch disease, and the wage problem}

A successful tourism sector earning foreign exchange will, under
Chapter~\ref{ch:money}'s dollarised regime, raise domestic prices and wages in
the sectors that serve it and squeeze everything that competes with imports ---
which is agriculture, which is the sector Chapter~\ref{ch:land} needs to grow.
This is the standard resource-boom pathology and it is aggravated here by the
labour constraint: there is no reserve of workers to expand into the boom
sector, so it draws them from the others.

Cuba has already run this experiment. In the 1990s the tourism wage drew doctors
into taxi driving and surgeons into renting rooms, which is where the inverted
pyramid of Chapter~\ref{ch:dollarreform} came from.

\begin{design}{Not repeating the 1990s}
\label{d:dutch}
\begin{rules}
\rl{1}{Public pay is the control} Design~\ref{d:pay} binds here. Public
professional pay is benchmarked against the tourism and private wage and moves
with it. A tourism strategy adopted without the corresponding public pay
settlement is a strategy for emptying the hospitals, and this book will not
propose one.

\rl{2}{Revenue to the fund} The accommodation levy and concession revenue are
rent under Design~\ref{d:rentrule} and go to the stabilisation and investment
fund, not to the recurrent budget --- which both dampens the domestic demand
impulse and stops a tourism boom from financing recurrent commitments that a
bust cannot sustain.

\rl{3}{Linkage over expansion} Policy effort goes into raising the domestic
content of each visitor's spending --- food, furniture, textiles, crafts, guiding
--- rather than into raising visitor numbers. Linkage is the instrument that
turns a boom sector into a diversified one, and it is the reason
Design~\ref{d:visitor}.4 is written as a tax differential rather than as an
exhortation.

\rl{4}{Regional weighting} Tourism development incentives are weighted toward
the eastern provinces, per Design~\ref{d:municipal}.2. Left alone, the industry
concentrates in Havana, Varadero and the northern cays, and the equalisation
formula then has to do all the work.
\end{rules}
\end{design}
